\documentclass{article}
\usepackage[T1]{fontenc}
\usepackage{hyperref}
\begin{document}
\section{Semantic source wrapping}
This short sentence stays together. A longer sentence contains a comma, then a semicolon; then a colon: and enough additional prose to exercise the preferred boundaries repeatedly without changing the printed paragraph.

An example uses e.g. an abbreviation, i.e. ordinary punctuation within prose, together with the decimal 3.14159 and the name \LaTeX{} in a sentence. Is the next sentence a question? Yes, it is!

This sentence uses a standalone -- dash, followed by a standalone --- longer dash, and some well-known hyphenated words that should remain intact.

\begin{abstract}
Short sentence. A longer sentence inside an indented environment has punctuation, together with enough words to check continuation widths and indentation.
\end{abstract}

A footnote\footnote{This is a footnote. Its macro argument is deliberately atomic, including the punctuation.} belongs to this sentence, which continues after its marker with more ordinary text.

A long URL stays indivisible: \url{https://example.org/a/very/long/path/with/commas,and;semicolons?value=one:two} is followed by a final sentence.

This final paragraph deliberately has no sentence-ending punctuation
\end{document}
